\documentclass[11pt,a4paper]{article}

\usepackage[margin=2.5cm]{geometry}
\usepackage{amsmath,amssymb,bm}
\usepackage[T1]{fontenc}
\usepackage{lmodern}

\title{Cubic Polynomial Trajectory Planning}
\author{}
\date{}

\begin{document}

\maketitle

\section{Cubic trajectory}

Consider a cubic polynomial trajectory defined over the time interval
$t\in[t_s,t_f]$:
\begin{equation}
    q(t)=a_3t^3+a_2t^2+a_1t+a_0.
\end{equation}
Its velocity and acceleration are
\begin{align}
    \dot q(t) &= 3a_3t^2+2a_2t+a_1,\\
    \ddot q(t) &= 6a_3t+2a_2.
\end{align}

The four polynomial coefficients are determined by prescribing the initial
and final positions and velocities:
\begin{align}
    q(t_s) &= q_s, & \dot q(t_s) &= \dot q_s,\\
    q(t_f) &= q_f, & \dot q(t_f) &= \dot q_f.
\end{align}

Substitution into the cubic polynomial gives the four boundary-condition
equations
\begin{align}
    a_3t_s^3+a_2t_s^2+a_1t_s+a_0 &= q_s,\\
    a_3t_f^3+a_2t_f^2+a_1t_f+a_0 &= q_f,\\
    3a_3t_s^2+2a_2t_s+a_1 &= \dot q_s,\\
    3a_3t_f^2+2a_2t_f+a_1 &= \dot q_f.
\end{align}

In matrix form, these equations become
\begin{equation}
\underbrace{\begin{bmatrix}
    t_s^3  & t_s^2 & t_s & 1\\
    t_f^3  & t_f^2 & t_f & 1\\
    3t_s^2 & 2t_s  & 1   & 0\\
    3t_f^2 & 2t_f  & 1   & 0
\end{bmatrix}}_{\bm{A}}
\underbrace{\begin{bmatrix}
    a_3\\a_2\\a_1\\a_0
\end{bmatrix}}_{\bm{a}}
=
\underbrace{\begin{bmatrix}
    q_s\\q_f\\\dot q_s\\\dot q_f
\end{bmatrix}}_{\bm{b}}.
\end{equation}
Therefore, the polynomial coefficients can be computed as
\begin{equation}
    \boxed{\bm{a}=\bm{A}^{-1}\bm{b}}.
\end{equation}
In numerical implementations, the linear system $\bm{A}\bm{a}=\bm{b}$
should normally be solved directly instead of explicitly forming
$\bm{A}^{-1}$.

\section{Solution using shifted time}

The closed-form solution is simpler if time is measured from the beginning
of the motion. Define
\begin{equation}
    \tau=t-t_s,
    \qquad
    T=t_f-t_s,
\end{equation}
so that $\tau\in[0,T]$. Write the trajectory as
\begin{equation}
    q(\tau)=a_0+a_1\tau+a_2\tau^2+a_3\tau^3.
\end{equation}
The velocity and acceleration are
\begin{align}
    \dot q(\tau) &= a_1+2a_2\tau+3a_3\tau^2,\\
    \ddot q(\tau) &= 2a_2+6a_3\tau.
\end{align}

The boundary conditions now read
\begin{align}
    q(0) &= q_s, & \dot q(0) &= \dot q_s,\\
    q(T) &= q_f, & \dot q(T) &= \dot q_f.
\end{align}
After substitution, one obtains
\begin{align}
    a_0 &= q_s,\\
    a_1 &= \dot q_s,\\
    a_0+a_1T+a_2T^2+a_3T^3 &= q_f,\\
    a_1+2a_2T+3a_3T^2 &= \dot q_f.
\end{align}

Equivalently,
\begin{equation}
\begin{bmatrix}
    1 & 0 & 0   & 0\\
    0 & 1 & 0   & 0\\
    1 & T & T^2 & T^3\\
    0 & 1 & 2T  & 3T^2
\end{bmatrix}
\begin{bmatrix}
    a_0\\a_1\\a_2\\a_3
\end{bmatrix}
=
\begin{bmatrix}
    q_s\\\dot q_s\\q_f\\\dot q_f
\end{bmatrix}.
\end{equation}

Solving this system gives
\begin{align}
    a_0 &= q_s,\\
    a_1 &= \dot q_s,\\
    a_2 &= \frac{3(q_f-q_s)}{T^2}
           -\frac{2\dot q_s+\dot q_f}{T},\\
    a_3 &= -\frac{2(q_f-q_s)}{T^3}
           +\frac{\dot q_s+\dot q_f}{T^2}.
\end{align}

Consequently, the complete trajectory is
\begin{equation}
\boxed{
\begin{aligned}
q(t)={}&q_s+\dot q_s(t-t_s)\\
&+\left[
\frac{3(q_f-q_s)}{T^2}
-\frac{2\dot q_s+\dot q_f}{T}
\right](t-t_s)^2\\
&+\left[
-\frac{2(q_f-q_s)}{T^3}
+\frac{\dot q_s+\dot q_f}{T^2}
\right](t-t_s)^3,
\end{aligned}}
\qquad T=t_f-t_s.
\end{equation}

\section{Zero initial and final velocities}

If the trajectory must start and finish at rest,
\begin{equation}
    \dot q_s=\dot q_f=0,
\end{equation}
the coefficients reduce to
\begin{align}
    a_0 &= q_s,
    &a_1 &= 0,
    &a_2 &= \frac{3(q_f-q_s)}{T^2},
    &a_3 &= -\frac{2(q_f-q_s)}{T^3}.
\end{align}
The corresponding trajectory is
\begin{equation}
\boxed{
q(t)=q_s
+3(q_f-q_s)\left(\frac{t-t_s}{T}\right)^2
-2(q_f-q_s)\left(\frac{t-t_s}{T}\right)^3
}.
\end{equation}

Introducing the normalized time
\begin{equation}
    s=\frac{t-t_s}{T},
    \qquad s\in[0,1],
\end{equation}
the same result can be written compactly as
\begin{equation}
    \boxed{q(s)=q_s+(q_f-q_s)(3s^2-2s^3)}.
\end{equation}

\end{document}
